\documentclass{article}
\usepackage{graphicx} % Required for inserting images
\usepackage{amsfonts}
\usepackage{amsmath}

\title{Lecture 4/8/26}
\author{}
\date{}

\begin{document}

\maketitle

\section{Sampling}
\begin{itemize}
    \item $X^{(i)}$ - sample one item from the population
    \item $X^{(1)}$ follows a sampling distribution from the population
    \item Sample: $\{X^{(1)} \dots X^{(N)}\}$
\end{itemize}

\section{Population Parameters and Test Statistic}
\begin{itemize}
    \item $\mu^*$ - population \textbf{parameter}, an attribute of the population
    \item $\mu$ - sample \textbf{statistic}, an attribute of the statistic used to estimate the population parameter
    \item $\mu = \mathbb{E}[X^{(1)} + \dots + X^{(N)}]$ is random, $\mu^*$ is not random
\end{itemize}

\section{Bias-Variance Decomposition}
\begin{itemize}
    \item $\mathbb{E}[(\mu^* - \mu)^2] = \underbrace{(\mathbb{E}[\mu^*] - \mu)^2}_{\text{Bias}^2} + \underbrace{\mathbb{E}[(\mu^* - E[\mu^*])^2]}_{\text{Var}(\mu^*)}$
    \item Bias - Accuracy of point estimate $\mu$ to parameter $\mu^*$
    \item Variance - How spread out we can expect statistic $\mu$ to be
    \item Optimal statistic: Low bias, low variance
    \item Consider the following example:
    \item Let $\mu^* = 0.6$
    \begin{align*}
        A &= 1.0 \\
        B &= 1.0 \\
        C &= 1.0 \\
        D &= 0.0 \\
        E &= 1.0
    \end{align*}
    Take all SRS of size 2, find $\mu$:
    \begin{align*}
        \{A, B\} = 1.0 \\
        \{A, C\} = 1.0 \\
        \{A, D\} = 0.5 \\
        \{A, E\} = 0.5 \\
        \{B, C\} = 1.0 \\
        \{B, D\} = 0.5 \\
        \{D, E\} = 0.5 \\
        \{C, D\} = 0.5 \\
        \{C, E\} = 0.5 \\
        \{D, E\} = 0.0
    \end{align*}
    \item $\mathbb{E}[\mu] = 0.6$. So, $\mu$ is an unbiased estimator for $\mu^*$.
    \item $\operatorname{Var}(\mu)=\operatorname{Var}\!\left(\frac{1}{n}\sum_{i=1}^{n}X^{(i)}\right)\overset{X^{(i)}\text{s independent}}{=}\frac{1}{n^2}\sum_{i=1}^{n}\operatorname{Var}(X^{(i)})\overset{X^{(i)}\text{s identical}}{=}\frac{\operatorname{Var}(X^{(i)})}{n}$
    \item So, for SRS, the bias term is 0, and the expected error is just the variance in $\mu$ via bias-variance decomposition.
    \item If the expected error is less than $\text{Var}(\mu)$, then there is \textbf{confirmation bias}.
\end{itemize}
\end{document}
